\documentclass[11pt, letterpaper]{article}

\usepackage[margin=1in]{geometry}
\usepackage{graphicx}
\usepackage{booktabs}
\usepackage{amsmath}
\usepackage{hyperref}
\usepackage{xcolor}
\usepackage{enumitem}
\usepackage{tabularx}

\graphicspath{{results/reports/}}

\definecolor{identical}{RGB}{34, 139, 34}
\definecolor{different}{RGB}{200, 30, 30}

\title{\textbf{Does Optimization Parallel Inference Change LLM Inference Outputs?}\\[0.4em]
\large A Concise Empirical Report on vLLM Configuration Effects}
\author{CSE 5429}
\date{\today}

\begin{document}

\maketitle

\section*{TL;DR}

We test whether vLLM system-level inference optimizations change LLM outputs
when the model, prompts, seed, and greedy decoding setup are fixed.  Among the
nine configurations that completed, six match the baseline exactly on all 20
prompts.  Two configurations change outputs: CUDA graphs match only 25\% of
baseline outputs, and batch processing matches 50\%.  The baseline determinism
check is perfectly repeatable, so the observed differences are tied to execution
configuration rather than ordinary run-to-run randomness.

\section{Question and Motivation}

Modern LLM serving systems improve throughput and latency with optimizations
such as chunked prefill, CUDA graphs, prefix caching, batching, and CPU
offloading.  These optimizations are usually treated as implementation details:
they should make inference faster, not change the generated text.  This report
tests that assumption directly.

The core question is simple: \emph{given the same model and prompt, do different
inference configurations produce the same output?}  This matters because output
reproducibility affects evaluation, debugging, and safety analysis.  If a model
answer changes when the serving path changes, then inference configuration is
part of the experimental condition, not merely an engineering detail.

\section{Experimental Setup}

The experiment uses \texttt{Qwen/Qwen2-1.5B-Instruct} served with vLLM v0.20.1
on a single NVIDIA RTX 4090.  All completed runs use greedy decoding with
temperature $=0$, seed $=42$, maximum context length 2048, and maximum generation
length 128.  The prompt set contains 20 prompts spanning short factual, medium
creative, long reasoning, code-generation, edge-case, and repetition-sensitive
inputs.  Each inference configuration is run in an isolated subprocess to avoid
GPU-memory state carrying over between configurations.

Table~\ref{tab:configs} summarizes the tested configurations.  C7 is optional
and was skipped in the final run because FP8 KV cache initialization failed in
this environment.

\begin{table}[h]
\centering
\caption{Inference configurations.}
\label{tab:configs}
\small
\begin{tabularx}{\textwidth}{@{} c l X @{} }
\toprule
\textbf{ID} & \textbf{Name} & \textbf{Purpose} \\
\midrule
C0 & Baseline & Reference configuration with eager execution \\
C1 & Chunked Prefill (512) & Tests small prefill chunking \\
C2 & Chunked Prefill (2048) & Tests larger prefill chunking \\
C3 & CUDA Graphs & Tests CUDA graph execution rather than eager mode \\
C4 & Prefix Caching & Tests KV cache prefix reuse \\
C5 & Batch Processing & Tests batched prompt scheduling \\
C6 & Low GPU Memory & Tests lower GPU memory utilization (40\%) \\
C7 & FP8 KV Cache & Optional cache-precision test; skipped \\
C8 & CPU Offload & Tests 4 GB CPU offload \\
C9 & Determinism & Repeats baseline three times \\
\bottomrule
\end{tabularx}
\end{table}

Outputs are compared with exact match, token match ratio, BLEU, ROUGE-L, edit
distance, and first divergent token position.  Exact match is the main metric:
it directly answers whether two generations are character-for-character
identical.

\section{Results}

Nine of ten configurations completed.  C7 was skipped after FP8 KV cache
initialization failed; all other configurations produced 20 outputs, except C9,
which produced 60 outputs because it repeats the baseline three times.
Table~\ref{tab:main-results} reports each completed non-baseline configuration
against C0.

\begin{table}[h]
\centering
\caption{Comparison against C0 baseline, averaged over 20 prompts.}
\label{tab:main-results}
\small
\begin{tabular}{@{} l c c c c @{} }
\toprule
\textbf{Config} & \textbf{Exact Match} & \textbf{BLEU} & \textbf{Edit Dist.} & \textbf{First Div.} \\
\midrule
C1 Chunked (512)  & \textcolor{identical}{1.000} & \textcolor{identical}{100.0} & \textcolor{identical}{0.0} & -- \\
C2 Chunked (2048) & \textcolor{identical}{1.000} & \textcolor{identical}{100.0} & \textcolor{identical}{0.0} & -- \\
C3 CUDA Graphs    & \textcolor{different}{0.250} & \textcolor{different}{69.7} & \textcolor{different}{154.2} & 28.0 \\
C4 Prefix Cache   & \textcolor{identical}{1.000} & \textcolor{identical}{100.0} & \textcolor{identical}{0.0} & -- \\
C5 Batch Proc.    & \textcolor{different}{0.500} & \textcolor{different}{78.7} & \textcolor{different}{111.5} & 16.7 \\
C6 Low Memory     & \textcolor{identical}{1.000} & \textcolor{identical}{100.0} & \textcolor{identical}{0.0} & -- \\
C8 CPU Offload    & \textcolor{identical}{1.000} & \textcolor{identical}{100.0} & \textcolor{identical}{0.0} & -- \\
C9 Determinism    & \textcolor{identical}{1.000} & \textcolor{identical}{100.0} & \textcolor{identical}{0.0} & -- \\
\bottomrule
\end{tabular}
\end{table}

The main result is sharp: chunked prefill, prefix caching, low memory, and CPU
offload are identical to the baseline; CUDA graphs and batch processing are not.
CUDA graphs match only 5 of 20 prompts exactly.  Batch processing matches 10 of
20 prompts exactly.  C9 confirms that the baseline is stable across repeated
runs, so the C3 and C5 differences are not explained by ordinary
nondeterministic reruns of C0.

\begin{figure}[h]
\centering
\includegraphics[width=0.72\textwidth]{heatmap_exact_match.png}
\caption{Pairwise exact match heatmap across completed configurations.}
\label{fig:exact-heatmap}
\end{figure}

\section{Discussion}

The results suggest a useful distinction between \emph{memory-management}
optimizations and \emph{execution-path} optimizations.  The memory-oriented
changes tested here---prefix caching, lower GPU memory utilization, and CPU
offload---do not change outputs.  They appear transparent for this model,
prompt set, and decoding setup.

CUDA graphs and batch processing are different.  Both change how GPU work is
scheduled or replayed.  A plausible explanation is numerical sensitivity:
different kernel execution paths can introduce small floating-point differences,
and greedy decoding can amplify a small logit-order change into a different
token choice.  This explanation is consistent with the observations, especially
the nonzero first-divergence positions, but it is not directly proven by these
experiments because we do not log intermediate logits or kernel-level numeric
differences.

The practical implication is that reproducibility-sensitive evaluations should
record the inference configuration, not only the model and prompt.  In this
experiment, the same model and greedy decoding setup can produce different text
when CUDA graphs or batching are enabled.  For throughput-oriented production
systems this may be acceptable; for exact reproducibility it is a real
experimental variable.

\section{Limitations}

This is a focused single-system experiment.  The conclusions are based on one
model, one GPU, one vLLM version, greedy decoding, and 20 prompts.  The FP8 KV
cache condition did not complete in this environment.  Larger models, multi-GPU
serving, sampling-based decoding, and direct logit comparison would be useful
follow-up experiments.

\section{Conclusion}

Most tested vLLM optimizations preserve exact outputs under this setup, but not
all do.  CUDA graphs and batch processing changed generated text even with a
fixed seed and temperature $=0$, while the baseline itself was perfectly
repeatable.  The central insight is therefore not that LLM inference is always
unstable, but that exact output reproducibility depends on the execution path.

\end{document}
